\subsection{3.4 有益性与安全攀登}\label{helpfulness-and-safety-climb}

有益性与安全（Helpfulness and Safety）RL 攀登以人类偏好、指令遵循、
可引导性（steerability）、安全、诚实性与风格为判据，优化 MAI-Thinking-1
的通用有益性。(Ouyang et al., 2022; Radford
et al., 2019; Li et al., 2024a) 与其他攀登阶段不同，本阶段聚焦于性能无法被
客观定义、也无法由机器验证的任务。

本节的组织方式是：先介绍 RL 攀登中使用的全部奖励信号
（第 3.4.1 节），随后介绍各领域的具体数据配方与奖励设计。

\subsection{3.4.1 奖励}\label{rewards}

与其他攀登阶段相比，有益性与安全攀登组合了更多样化的奖励类型，
以引导模型行为中较为主观的方面。我们结合了基于人类偏好数据训练的
奖励模型 (Lambert et al., 2025)、AI 裁判反馈（通常由评分准则引导）(Gu et al., 2024)，
以及额外的可验证奖励
(Peng et al., 2025; Pyatkin et al., 2026)，形成聚合奖励
信号。(Lambert et al., 2024; Wang et al., 2024c)

奖励模型。我们的奖励模型基于 MAI-Base-1 的后训练版本，
我们对其微调，以预测以文本 token 表达的人类偏好。
该模型仅使用由多家供应商的人类标注者收集的人类偏好数据
进行训练。我们采用 Liu et al.~(2025c) 所述的一套稳健的奖励劫持
（reward hacking）缓解管线，
以应对训练数据中的相关偏差。

训练。对于上下文 c 和 k 路并排比较（k-way side-by-side）的响应
\(y _ { 1 } , ~ . . . , ~ y _ { k } ,\) ，以及对应分数
\(s _ { 1 } , \ldots , s _ { k } \in [ 1 ; 5 ]\) ，奖励模型的输入为

\[
c \texttt{<|im\_sep|>} y_1 \texttt{<|im\_sep|>} y_2 \texttt{<|im\_sep|>} \ldots \texttt{<|im\_sep|>} y_k \texttt{<|im\_sep|>}
\]

其中训练目标是序列
\(s _ { 1 } \ldots s _ { k } ,\) ，通过 SFT 训练。

推理。将多个候选响应放入同一个奖励模型上下文，可以在不同响应之间
对逐点（pointwise）质量分数进行更好的校准。然而，由于奖励模型
的自回归特性，这也增加了分数
\(y _ { i , i > 1 }\) 的噪声。为应对这一点，我们循环地应用奖励模型：
对于给定上下文 c 和 k 个响应
\(y _ { 1 } , \ldots , y _ { k }\) ，我们按响应的循环置换
\(( y _ { 1 } , \dots , y _ { k } ) , ( y _ { 2 } , \dots , y _ { k } , y _ { 1 } ) , \dots , ( y _ { k } , y _ { 1 } , \dots , y _ { k - 1 } )\)
对奖励模型提示 k 次
。对于这 k 次推理调用中的每一次，我们只解码第一个 token，
并查看该 token 的完整概率分布；第 i 次调用对应的即
第 i 个候选响应的奖励分数 \(s _ { i }\) 。奖励信号
\(R _ { \mathrm { R M } } ( c , y _ { i } )\) 随后被设置为
\(y _ { i }\) 被评为最高质量
\(( s _ { i } = 5 )\) 的概率。我们发现，与我们考虑过的其他
公式化方案相比，这能提供更稳定的攀登信号。

评估。我们结合两方面的数据评估奖励模型：由标注者评分的实际
训练 rollout，以及奖励模型训练数据的验证集划分。

AI 裁判。除了微调奖励模型提供的人类偏好信号外，我们还采用
AI 裁判 (Gu et al., 2024; Tan et
al., 2025) 提供反馈，其优势是能够快速适配并针对
任何给定上下文定制。特别是，AI 反馈提供了一种快速杠杆，
可以有针对性地塑造模型行为，而无需承担
奖励模型数据收集和重新训练的时间成本。

可验证奖励。我们采用可验证奖励来提升某些领域的能力，例如指令遵循——
这类领域中，对约束的遵循情况可以被直接检查。例如，对于要求
``respond in a single paragraph''（用单段话回答）或 ``use fewer
than 10 words''（用不超过 10 个词）的上下文，我们会纳入可验证奖励。
与不可验证奖励相比，我们发现可验证奖励
更不易被奖励劫持 (Zhao et al., 2025b; Wen et
al., 2024)，对多轮训练（multi-epoching）的敏感性更低，并且通常有助于
稳定训练。

可验证奖励还被用于缓解不可验证奖励中的偏差。例如，AI 反馈评分准则
往往会对长度和风格要素产生向上压力，我们同时使用
奖励模型和可验证奖励来缓解这一点。具体到长度，我们
离线（offline）确定每个上下文可接受的响应长度分布，并对落在
预定义分位数区间之外的响应施加惩罚。这
使响应长度在训练过程中保持稳定，既不过度约束响应，
也不会引入一种可被劫持的奖励——即通过把响应压缩得
过分简短来钻空子的奖励。

奖励组合。在为有益性与安全攀登优化这些奖励时，
会出现两个挑战。(Wang et al., 2024c) 首先，
不同类型的奖励处于不同的尺度，无法直接
比较。其次，奖励分布本身依赖上下文：
对某些提示词而言窄，对另一些则宽。天真地将这些奖励相加
会导致幅值最大的信号——无论其重要性如何——占据主导。
此外，虽然我们希望联合优化
各项奖励，但某些标准是不可让步的：例如，
写得很好的响应如果不安全，仍然不可接受，无论
其质量如何。我们通过两种互补的
策略应对这些挑战，并根据上下文有选择地应用。

词典序奖励塑形（Lexicographic reward shaping）。对于一组上下文，当
一组 rollout 在更高重要性的奖励上得分完全相同时，较低重要性的
奖励才被激活。这会带来严格的优先级排序：
只有在 rollout 组内主要奖励打成平手时，次级奖励才会影响
梯度 (Skalse et al., 2022)。由于
它基于组内相对比较，这一公式化方案
也对每个奖励的绝对尺度不敏感。

门控奖励应用（Gated reward application）。对于其他上下文，更高重要性的奖励
必须先达到最低性能水平，之后才会应用
较低重要性的奖励 (Achiam et al., 2017)。安全就是
典型例子：不安全的响应获得最低奖励，并且
永远不会被按响应质量评分。



这两种策略都以条件式或基于排序（rank-based）的逻辑取代加性组合，
从而规避跨尺度校准问题，
确保高优先级的目标绝不会为了
低优先级目标上的收益而被牺牲。

\subsection{3.4.2 指令遵循与可引导性}\label{instruction-following-and-steerability}

指令遵循（instruction following，IF）是 LLM 的核心能力：模型应遵循
对话中用户直接指定的约束、API 用户通过开发者指令指定的约束，
以及平台所有者通过特权系统指令指定的约束。模型应
能按照预定义的优先级，在这些格式与行为
约束之间被引导 (Wallace et al., 2024)。为构建稳健的指令遵循
能力，我们构建了一个覆盖多种
约束、场景与复杂度层级的数据集，数据来源包括
合成数据与专家人工标注。

数据。我们使用两个不同的数据来源：专家编写的上下文和
合成数据。我们发现，专家编写提示词中的复杂
约束有助于启动能力，而合成数据则能实现
最大覆盖。

合成数据生成。受先前工作启发 (Xu et al., 2025a;
Dong et al., 2025)，我们构建了一条灵活的多阶段管线，
生成带评估准则的真实指令与场景。
它首先利用 LLM 生成提示指令和模型规格，
该 LLM 由一个手工整理的约束分类体系（详见附录 F）引导，
并配有一组多样化的种子。随后生成
场景：这些场景是多语言的，覆盖简短与长篇对话，
包含系统、开发者和用户消息等情形，
并跨越超过 40 个领域，包括编程、写作、
分析与旅行规划。此外，
我们还加入带有冲突的系统、开发者和
用户指令的对抗性案例，以训练模型
尊重指令层级（instruction hierarchy）。最后的批判与改写步骤
会基于自然度、准则对齐度和
接地性（groundedness）评估并更新每个场景。

数据过滤。我们通过多轮过滤来保持高质量数据，
综合了各种质量启发式规则、复杂度过滤器以及
拒绝采样。特别是，评分准则需要经过验证，以确保
自包含性、无歧义性以及与所声明约束的对齐，
同时提示词需要经过我们的安全政策筛查。
参照 An et al.~(2025) 的做法，我们还通过通过率（pass-rate）分析控制难度水平。
我们只保留具有挑战性但
能被我们的模型解决的样本。

奖励设计。为防止 RL 训练中出现退化响应和奖励劫持，
我们使用混合奖励信号。对于具有确定性验证器的约束，
我们使用基于规则的检查。这提供了快速、确定性、
且校准良好的奖励信号。

此外，一个 LLM 裁判（第 3.4.1 节）会独立评估原子化
评分准则，为每条准则产生一个二元判断。多次
判断结果取平均以增强稳健性。最后，我们训练好的奖励模型
评估一般响应质量，对针对约束的
IF 裁判形成补充。随后，奖励按照第 3.4.1 节中的
词典序评分聚合方案组合，其中 IF 专属奖励作为
主要信号。

\subsection{3.4.3 安全}\label{safety}

与我们对支持人类自主性的承诺一致，我们将
安全定义为模型在提供有益响应的同时
保持符合我们政策的能力 (Microsoft, 2022)。

数据。为使训练与这一目标保持一致，我们为数据整理
开发了一套分类体系。安全数据针对两种互补的失败
模式：不安全合规（unsafe compliance），即模型完成了
本应拒绝的请求；以及过度拒绝（over-refusal），即模型不必要地
拒绝了合法请求。每个候选样本都依据政策分类体系
进行标注，并被归入两类提示词之一：



表 8. 有害提示词与边界提示词的来源。

{\def\LTcaptype{none} % do not increment counter
\begin{longtable}[]{@{}lll@{}}
\toprule\noalign{}
\endhead
\bottomrule\noalign{}
\endlastfoot
类别 & 管线 & 来源 \\
\multirow{2}{*}{有害} & 人类红队测试 & 供应商编写的提示词
内部红队演练 \\
& 自动化攻击 & 基于模板的攻击，如 PyRIT (Munoz et al.,
2024) 非交互式 LLM 生成，如 PAP (Zeng et al., 2024a)
交互式 LLM 驱动，如 TAP (Mehrotra et al., 2024) \\
边界 & Do-Not-Refuse 分片 能力数据 & 跨数据代际保留的精选数据
经安全管线标注进入 do-not-refuse 分片 \\
\end{longtable}
}

• 有害提示词：政策禁止提供部分或全部帮助的请求。
响应策略为完全拒绝或部分拒绝
（拒绝不安全的部分，同时提供安全的替代方案）。

• 边界提示词：涉及敏感领域但在政策范围内仍可回答的请求。
响应策略为 do-not-refuse（不拒绝）：
提供有边界、有帮助的回答，而不是含糊其辞或拒绝。

数据来源。表 8 总结了各类提示词的获取方式。
有害提示词通过人类红队测试
和自动化对抗生成两种途径收集。边界提示词来自
跨数据代际保留的现有 do-not-refuse 分片，
以及经由安全标注管线路由并被选入
do-not-refuse 分片的能力数据，从而让
模型接触到应当保持可回答的政策邻近型（policy-adjacent）请求。

奖励设计。每个模型响应由安全裁判沿三个维度评分：

• 政策合规（Policy compliance）：衡量响应是否违反我们的安全
政策。

• 响应参与度（Response engagement）：将响应的参与程度与
预期目标（拒绝、部分拒绝或 do-not-refuse）进行比较，
对过度拒绝以及严重请求上的合规行为都施加惩罚。

• 响应风格（Response style）：衡量响应是否遵循预期的
语气和指导原则，例如对敏感的自残请求，
要承认其困难性，而不是说教或
向用户复述创伤事件。

裁判分数与来自奖励模型的奖励信号、
（视数据来源而定）额外的 AI 裁判信号相结合。
简单的加权平均是不够的：标量奖励
对政策不合规的响应仍可能保持为正。\footnote{在一项成对 rollout 审计中，期望安全 Likert 分数与政策合规字段只有中等一致性（Pearson 0.293，Spearman 0.344），且 87.8\% 的政策不合规响应获得的奖励模型分数不低于 3。}
因此我们使用安全门控聚合器：按第 3.4.1 节所述，
政策合规字段对奖励进行门控。如果安全裁判判定
响应不合规，则该 rollout 获得最低奖励，无论
其他分数如何；否则使用正常的加权混合。

\subsection{3.4.4 诚实性}\label{honesty}

我们将诚实性定义为：模型在知道答案时
能给出正确响应的能力，以及在不知道答案时恰当表达不确定的能力。
同时，模型不应过度表达不确定，因为
单纯最小化事实错误只会促使模型做出更少的论断
(Min et al., 2023)。因此，目标必须
在事实精确性与信息量之间取得平衡。



表 9. 我们风格指南中目标行为描述的示例。

{\def\LTcaptype{none} % do not increment counter
\begin{longtable}[]{@{}|l|l|@{}}
\toprule\noalign{}
\endhead
\bottomrule\noalign{}
\endlastfoot
\hline
关注领域 & 示例指导 \\
\hline
表情符号 & 尽量少用——仅在明确合适的休闲语境中使用。
不要用表情符号作为列表标题或署名落款。 \\
\hline
表格与 Markdown & 项目符号列表和表格用于罗列
多个要点、给出示例或呈现数据。当回答
确需多个不同部分时，用标题将其拆分。 \\
\hline
语气与正式程度 & 默认采用专业语气，仅在用户直接
要求时才采用轻松或富有感染力的语气。 \\
\hline
版面布局 & 尽量少用开场白或收尾套话，不写谄媚式的
引言。信息层级清晰。回答
便于快速浏览，重要信息放在前面。 \\
\hline
密度 & 避免大段文字堆砌。粗体字要有策略地使用。
段落之间留出间距。句子保持足够
简短易读。 \\
\hline
\end{longtable}
}

数据。我们汇集了多样化的手工整理数据，来源包括供应商、
经 PII 过滤的消费者 Copilot 日志\footnote{消费者 Copilot 日志数据仅面向未选择退出将其数据用于模型训练的用户收集，并排除微软 Copilot 隐私常见问题（Privacy FAQ for Microsoft Copilot \textbar{} Microsoft Support）中所列的部分用户。}以及合成生成的数据。
受先前工作启发 (Zhang et al., 2024; Wei et al., 2024b,c)，
我们的数据覆盖一个难度谱系：(a) 既有事实性查询，
包括短格式与长格式，其响应可以与
参考答案进行精确核验；(b) 有挑战性的事实性查询，
面向长尾或冷门主题——这些主题上模型的覆盖
预期不稳定，参考标签通过检索增强式核验生成；
以及 (c) 虚假前提查询，即问题中包含
错误的预设（presupposition），且不存在正确的肯定性回答。
覆盖这些边界情形促使模型保持事实
完整性，并且只在缺乏相关知识时才表达不确定。

奖励设计。在事实性评分中，每个 RL 样本的参考标签均
通过检索增强生成与核验离线（offline）生成。
随后每个模型响应由 LLM 裁判沿
两个维度——事实性与置信度——评分，得到五类之一：
CONFIDENT\_CORRECT（自信且正确）、UNCONFIDENT\_CORRECT（不自信但正确）、NOT\_ATTEMPTED
（未作答）、UNCONFIDENT\_INCORRECT（不自信且错误）和 CONFIDENT\_INCORRECT（自信但错误）。
这些类别通过加权求和
合并为标量奖励：自信且正确的响应
获得最高奖励，自信的幻觉受到最严厉的
惩罚，放弃作答获得中性奖励，而不自信但正确的
响应获得降低的奖励，以抑制过度表达不确定。

\subsection{3.4.5 风格}\label{style}

我们制定了一份风格指南（style guide），规定优质输出应是什么样：亲切而
不谄媚、结构易于扫读、语气依语境校准
而非简单模仿用户。风格指南还涵盖何时以及
如何使用表情符号、正式程度、数学与代码的记法，以及
总体文本密度。表 9 给出了我们训练模型遵循的
指导原则的高层示例。

数据。风格评分基于广泛的数据集：经 PII 过滤的微软消费者
Copilot 日志、供应商编写的上下文（静态与交互式），以及
Arena 对话。这些数据覆盖低至中等难度的提示词，
排除复杂的指令遵循、编程、数学和 STEM。
提示词按用户意图分类（如创意写作、实用
指导、信息寻求和闲聊），并主动采集
以针对模型较弱的领域。



表 10. 合并 SFT 数据混合按能力的组成。

{\def\LTcaptype{none} % do not increment counter
\begin{longtable}[]{@{}|l|l|l|@{}}
\toprule\noalign{}
\endhead
\bottomrule\noalign{}
\endlastfoot
\hline
能力 & 样本权重 & Token 权重 \\
\hline
STEM 与编程 & 56\% & 89\% \\
\hline
智能体能力 & 11\% & 9\% \\
\hline
通用有益性与安全 & 33\% & 2\% \\
\hline
\end{longtable}
}

奖励设计。除了第 3.4.1 节介绍的奖励模型外，
我们还使用带提示词的 LLM 裁判沿特定轴为输出评分，
并对不良行为施加惩罚。这些带提示词的裁判在
0、1、2 的整数尺度上为响应评分，分别对应重大、轻微或
无风格问题。我们发现，粗粒度裁判
优于更细粒度的裁判或按提示词定制的评分准则，
因为它们让 LLM 裁判有更大的灵活性，
能针对给定的提示词与响应恰当地解释准则，
从而让模型更难劫持裁判。风格裁判
仅在可验证奖励与安全约束均已满足后才被应用，
并根据领域与其他裁判一同加权。
